

% ---------------------------------------------------------------------------
% A3b -- pushes the "mandatory-dominated, no wave" shape back to a PRE-SAMPLE
% (2016) baseline from an INDEPENDENT source (TRE-SP SAC-JE), the only 2016
% litigation snapshot we have. Compatibility is deliberately loose: 2016 zone
% IDs (SAC-JE, 1-427) do not align with the SIG panel's zone-codes (1-107), so
% this is a STATE-LEVEL composition comparison, never zone-by-zone; source
% changes across the divide, so only within-source SHARES are read (no count
% trend). Numbers come from lawsuit_composition_sp.tex, generated by
% 04_analysis/08_lawsuit_composition_sp.py -- nothing hard-coded.
% ---------------------------------------------------------------------------
\begin{frame}{Appendix --- SP Litigation Composition, 2016--2024}
\hypertarget{app:sp2016}{}
\footnotesize
\noindent A separate TRE-SP filing extract (S\~{a}o Paulo, all 645 municipalities) gives the only \textbf{pre-2020} snapshot of the docket. It carries the flat, mandatory-dominated shape back to before the sample period. The mandatory core (candidacy registration and campaign accounts) is the overwhelming majority of filings in every cycle ($\sim$95\%, shaded rows), and the adversarial slice is small and rises only gently across the three cycles (shaded, bottom).
\smallskip
\begin{columns}[T]
\begin{column}{0.50\linewidth}
\centering
\resizebox{\linewidth}{!}{% Auto-generated by code/04_analysis/08_lawsuit_composition_sp.py
% Do not edit manually -- rerun the generating script to update

\begin{tabular}{lccc}
\toprule\toprule
 & \textbf{2016} & \textbf{2020} & \textbf{2024} \\
 & \scriptsize{TRE-SP} & \scriptsize{SIG} & \scriptsize{SIG} \\
\midrule
\multicolumn{4}{l}{\textit{Mandatory / ancillary filings}} \\
\quad Candidacy registration & 62.8 & 52.2 & 49.9 \\
\quad Campaign accounts & 32.1 & 42.7 & 44.0 \\
\rowcolor{mylight}
\textbf{Mandatory core} & \textbf{94.9} & \textbf{94.9} & \textbf{93.8} \\
\midrule
\multicolumn{4}{l}{\textit{Adversarial filings}} \\
\quad Representa\c{c}\~{a}o & 2.45 & 2.58 & 2.97 \\
\quad Propaganda$^{\dagger}$ & 0.00 & 0.67 & 1.01 \\
\quad AIJE / AIME / criminal & 0.44 & 0.31 & 0.27 \\
\rowcolor{mylight}
\textbf{Adversarial} & \textbf{2.89} & \textbf{3.57} & \textbf{4.25} \\
\midrule
Administrative \& other & 2.2 & 1.5 & 1.9 \\
\bottomrule\bottomrule
\end{tabular}
}
\end{column}
\begin{column}{0.48\linewidth}
{\scriptsize
The 2016 baseline does two kinds of work. It confirms that the mandatory-filing exclusion is structural rather than a 2020-era SIG artifact, since the $\sim$95\% core predates our data. It also extends the ``expanded role, not volume'' reading to a pre-sample year: against a clean 2016 benchmark the adversarial arena still does not expand sharply.\par
\smallskip
$^{\dagger}$\,The propaganda row requires care. In 2016 there was no dedicated propaganda class, and those complaints filed as \emph{Representa\c{c}\~{a}o}, so the apparent $0\to1\%$ rise is partly TSE reclassification rather than behavior. The total adversarial share is comparable across cycles; its internal split is not.\par}
\end{column}
\end{columns}
\smallskip
{\centering\scriptsize Only some comparisons are licensed here. Zone identifiers do not align across sources (2016 SAC-JE spans zones 1--427; the SIG panel carries codes 1--107), so this is a state-level composition comparison rather than a zone-by-zone one. The source also changes at the 2016/2020 divide, so only within-source shares are read, never a raw-count trend.\par}
\smallskip
\centerline{\hyperlink{app:recomp}{\beamerreturnbutton{Recomposition (2020--24)}}\quad\hyperlink{src:motivation}{\beamerreturnbutton{Back}}}
\end{frame}
